\section{Discrete-Time SIR Parameter Estimation}

\subsection{Cramér-Rao Lower Bound: Theory and Application}

The \emph{Cramér-Rao Lower Bound (CRLB)} provides a fundamental limit on the precision of any unbiased estimator, establishing the minimum achievable variance for parameter estimation from noisy observations.

\subsubsection{General Theory}

Consider a parameter vector $\boldsymbol{\phi} \in \mathbb{R}^p$ to be estimated from observations $\mathbf{Y} \in \mathbb{R}^n$ with probability density $f(\mathbf{Y}|\boldsymbol{\phi})$. The \emph{Fisher Information Matrix} is defined as:
\begin{equation}
[\mathbf{J}]_{ij} = -\mathbb{E}\left[\frac{\partial^2 \log f(\mathbf{Y}|\boldsymbol{\phi})}{\partial \phi_i \partial \phi_j}\right] = \mathbb{E}\left[\frac{\partial \log f(\mathbf{Y}|\boldsymbol{\phi})}{\partial \phi_i}\frac{\partial \log f(\mathbf{Y}|\boldsymbol{\phi})}{\partial \phi_j}\right]
\end{equation}

\emph{CRLB Statement}: For any unbiased estimator $\hat{\boldsymbol{\phi}}$ of $\boldsymbol{\phi}$:
\begin{equation}
\text{Cov}(\hat{\boldsymbol{\phi}}) \succeq \mathbf{J}^{-1}
\end{equation}
where $\mathbf{A} \succeq \mathbf{B}$ means $\mathbf{A} - \mathbf{B}$ is positive semidefinite.

For Gaussian observations $\mathbf{Y} \sim \mathcal{N}(\boldsymbol{\mu}(\boldsymbol{\phi}), \sigma^2 \mathbf{I})$, the Fisher Information simplifies to:
\begin{equation}
\mathbf{J} = \frac{1}{\sigma^2} \mathbf{G}^T \mathbf{G}
\end{equation}
where $\mathbf{G} = \frac{\partial \boldsymbol{\mu}}{\partial \boldsymbol{\phi}}$ is the Jacobian of the mean function.

\subsubsection{Application to Epidemic Connectivity Inference}

In epidemic modeling, we seek to estimate \emph{connectivity parameters} $\theta$ that quantify transmission between regions. The fundamental question is: *Can inter-regional connectivity be reliably inferred from epidemic surveillance data?*

The CRLB provides the theoretical foundation to answer this question by establishing the \emph{information-theoretic limits} of connectivity estimation. If the CRLB shows that even optimal unbiased estimators have unacceptably large variance, then connectivity inference is fundamentally impossible regardless of estimation method.

We derive the Cramér-Rao lower bound for estimating connectivity
parameters from discrete-time epidemic observations using a
compartmental SIR model.

\subsection{Discrete-Time SIR Dynamics}

The compartments evolve according to the discrete-time SIR model:
\begin{align}\begin{split}\label{eq:sir}
S_i(t+1) &= S_i(t) - \beta(t) S_i(t) \sum_j C_{ij} I_j(t) \\
I_i(t+1) &= I_i(t) + \beta(t) S_i(t) \sum_j C_{ij} I_j(t) - \gamma I_i(t) \\
R_i(t+1) &= R_i(t) + \gamma I_i(t)
\end{split}\end{align}

For two regions with connectivity matrix:
\begin{equation}
\mathbf{C} = \begin{bmatrix} 1 & \theta \\ \theta & 1 \end{bmatrix}
\end{equation}

The dynamics become:
\begin{align}
S_1(t+1) &= S_1(t) - \beta(t) S_1(t)[I_1(t) + \theta I_2(t)] \\
S_2(t+1) &= S_2(t) - \beta(t) S_2(t)[\theta I_1(t) + I_2(t)] \\
I_1(t+1) &= I_1(t) + \beta(t) S_1(t)[I_1(t) + \theta I_2(t)] - \gamma I_1(t) \\
I_2(t+1) &= I_2(t) + \beta(t) S_2(t)[\theta I_1(t) + I_2(t)] - \gamma I_2(t)
\end{align}

with initial conditions $S_i(0), I_i(0)$ and $R_i(0) = 1 - S_i(0) - I_i(0)$.

\subsection{Problem Formulation}

Consider two regions observed over $T$ discrete time steps ($t = 0, 1, \ldots, T-1$). The parameter vector to be estimated is:
\begin{equation}
\boldsymbol{\phi} = [\theta, S_1(0), I_1(0), S_2(0), I_2(0)]^T \in \mathbb{R}^5
\end{equation}
where $\theta$ is the connectivity parameter and $S_i(0), I_i(0)$ are initial conditions.

The observation vector consists of new infections at each time step:
\begin{equation}
\mathbf{Y} = \begin{bmatrix} Y_1(0) \\ Y_2(0) \\ \vdots \\ Y_1(T-1) \\ Y_2(T-1) \end{bmatrix} \in \mathbb{R}^{2T}
\end{equation}

The observations are the new infections with additive noise:
\begin{equation}
\begin{bmatrix} Y_1(t) \\ Y_2(t) \end{bmatrix} = \beta(t) \begin{bmatrix} S_1(t)[I_1(t) + \theta I_2(t)] \\ S_2(t)[\theta I_1(t) + I_2(t)] \end{bmatrix} + \begin{bmatrix} \varepsilon_1(t) \\ \varepsilon_2(t) \end{bmatrix}
\end{equation}
where $\varepsilon_i(t) \sim \mathcal{N}(0, \sigma^2)$ independently, for $t = 0, 1, \ldots, T-1$.

\subsection{Mean Vector and Parameter Dependencies}

The mean vector is $\boldsymbol{\mu}(\boldsymbol{\phi}) =
\mathbb{E}[\mathbf{Y}|\boldsymbol{\phi}]$ where each component depends
on the discrete-time evolution:
\begin{equation}
\boldsymbol{\mu}(\boldsymbol{\phi}) = \begin{bmatrix} 
\mu_1(0) \\ \mu_2(0) \\ \vdots \\ \mu_1(T-1) \\ \mu_2(T-1) 
\end{bmatrix} = \begin{bmatrix}
\beta(0) S_1(0)[I_1(0) + \theta I_2(0)] \\
\beta(0) S_2(0)[\theta I_1(0) + I_2(0)] \\
%\beta(t) S_1(1)[I_1(1) + \theta I_2(1)] \\
%\beta(t) S_2(1)[\theta I_1(1) + I_2(1)] \\
\vdots \\
\beta(T-1) S_1(T-1)[I_1(T-1) + \theta I_2(T-1)] \\
\beta(T-1) S_2(T-1)[\theta I_1(T-1) + I_2(T-1)]
\end{bmatrix}
\end{equation} The Jacobian matrix is:

\begin{align}
  \mathbf{G} &= \frac{\partial \boldsymbol{\mu}}{\partial \boldsymbol{\phi}} =
  \begin{bmatrix}
    |& |& |& |& |\\
    \frac{\partial \boldsymbol{\mu}}{\partial \theta} & 
    \frac{\partial \boldsymbol{\mu}}{\partial S_1(0)} & 
    \frac{\partial \boldsymbol{\mu}}{\partial I_1(0)} & 
    \frac{\partial \boldsymbol{\mu}}{\partial S_2(0)} & 
    \frac{\partial \boldsymbol{\mu}}{\partial I_2(0)}
    \\|& |& |& |& |
  \end{bmatrix} \in \mathbb{R}^{2T \times 5}\\
  G_{kl} &= \frac{\partial \boldsymbol{\mu}_k}{\partial \boldsymbol{\phi}_l}
\end{align}
Define
\begin{equation}
  \Omega = \frac{\partial C}{\partial \theta} :=
  \begin{bmatrix} 0 & 1 \\ 1 & 0 \end{bmatrix}.
\end{equation}
Letting $y \in \{S_1(0), S_2(0), I_1(0), I_2(0)\}$, the partial
derivatives are:
\begin{align}
\frac{\partial \mu(t)}{\partial \theta} &= \beta(t) S(t) \circ \Omega I(t) \in \mathbb{R}^2 \\
\frac{\partial \mu(t)}{\partial y} &= \beta(t)\left [ \frac{\partial S(t)}{\partial y} \circ C I(t) + S(t) \circ C \frac{\partial I(t)}{\partial y}\right ] \in \mathbb{R}^2
\end{align}

We need to find derivatives with respect to initial conditions $y$, i.e.~$\frac{\partial S(t)}{\partial y}$ and $\frac{\partial I(t)}{\partial y}$. These evolve according:

\begin{align}\begin{split}
  \frac{\partial S(t+1)}{\partial y} &= \frac{\partial S(t)}{\partial y} - \beta(t) \left[ \frac{\partial S(t)}{\partial y} \circ  CI(t) + S(t)\circ C \frac{\partial I(t)}{\partial y} \right] \\
  & = \frac{\partial S(t)}{\partial y} - \frac{\partial \mu(t)}{\partial y} \\
  %
  %
  %
  \frac{\partial I(t+1)}{\partial y} &= (1-\gamma) \frac{\partial I(t)}{\partial y} + \frac{\partial \mu(t)}{\partial y}.
\end{split}\end{align}

%% \begin{align}
%% \frac{\partial S_i(t+1)}{\partial I_j(0)} &= \frac{\partial S_i(t)}{\partial I_j(0)} - \beta(t) \left[ \frac{\partial S_i(t)}{\partial I_j(0)} \sum_k C_{ik} I_k(t) + S_i(t) \sum_k C_{ik} \frac{\partial I_k(t)}{\partial I_j(0)} \right] \\
%% \frac{\partial I_i(t+1)}{\partial I_j(0)} &= \frac{\partial I_i(t)}{\partial I_j(0)} + \beta(t) \left[ \frac{\partial S_i(t)}{\partial I_j(0)} \sum_k C_{ik} I_k(t) + S_i(t) \sum_k C_{ik} \frac{\partial I_k(t)}{\partial I_j(0)} \right] - \gamma \frac{\partial I_i(t)}{\partial I_j(0)}
%% \end{align}

with initial conditions:
\begin{align}
\frac{\partial S_i(0)}{\partial S_j(0)} &= \delta_{ij}, \quad \frac{\partial I_i(0)}{\partial S_j(0)} = 0 \\
\frac{\partial S_i(0)}{\partial I_j(0)} &= 0, \quad \frac{\partial I_i(0)}{\partial I_j(0)} = \delta_{ij}
\end{align}

\subsubsection{Vectorized Sensitivity Equations}

The above scalar equations can be written compactly using matrix operations. Define the $2 \times 2$ sensitivity matrices:
\begin{align}
\frac{\partial \mathbf{S}(t)}{\partial \mathbf{S}(0)} &= \begin{bmatrix} \frac{\partial S_1(t)}{\partial S_1(0)} & \frac{\partial S_1(t)}{\partial S_2(0)} \\ \frac{\partial S_2(t)}{\partial S_1(0)} & \frac{\partial S_2(t)}{\partial S_2(0)} \end{bmatrix}, \quad \frac{\partial \mathbf{I}(t)}{\partial \mathbf{S}(0)} = \begin{bmatrix} \frac{\partial I_1(t)}{\partial S_1(0)} & \frac{\partial I_1(t)}{\partial S_2(0)} \\ \frac{\partial I_2(t)}{\partial S_1(0)} & \frac{\partial I_2(t)}{\partial S_2(0)} \end{bmatrix}
\end{align}

and similarly for derivatives with respect to $\mathbf{I}(0)$. Let $\mathbf{S}(t) = [S_1(t), S_2(t)]^T$, $\mathbf{I}(t) = [I_1(t), I_2(t)]^T$, and $\mathbf{C} = \begin{bmatrix} 1 & \theta \\ \theta & 1 \end{bmatrix}$.

The vectorized sensitivity equations are:

\begin{align}
  \frac{\partial \mathbf{S}(t+1)}{\partial \mathbf{S}(0)} &= &&\frac{\partial \mathbf{S}(t)}{\partial \mathbf{S}(0)} - \beta(t) \left[ \mathbf{C} \mathbf{I}(t) \circ \frac{\partial \mathbf{S}(t)}{\partial \mathbf{S}(0)} + \mathbf{S}(t) \circ \mathbf{C} \frac{\partial \mathbf{I}(t)}{\partial \mathbf{S}(0)} \right] \\
  %
  %
  %
  \frac{\partial \mathbf{I}(t+1)}{\partial \mathbf{S}(0)} &= (1-\gamma) &&\frac{\partial \mathbf{I}(t)}{\partial \mathbf{S}(0)} + \beta(t) \left[ \mathbf{C} \mathbf{I}(t) \circ \frac{\partial \mathbf{S}(t)}{\partial \mathbf{S}(0)} + \mathbf{S}(t) \circ \mathbf{C} \frac{\partial \mathbf{I}(t)}{\partial \mathbf{S}(0)} \right] \\
  %
  %
  %
  \frac{\partial \mathbf{S}(t+1)}{\partial \mathbf{I}(0)} &= &&\frac{\partial \mathbf{S}(t)}{\partial \mathbf{I}(0)} - \beta(t) \left[ \mathbf{C} \mathbf{I}(t) \circ \frac{\partial \mathbf{S}(t)}{\partial \mathbf{I}(0)} + \mathbf{S}(t) \circ \mathbf{C} \frac{\partial \mathbf{I}(t)}{\partial \mathbf{I}(0)} \right] \\
  %
  %
  %
  \frac{\partial \mathbf{I}(t+1)}{\partial \mathbf{I}(0)} &= (1-\gamma)&&\frac{\partial \mathbf{I}(t)}{\partial \mathbf{I}(0)} + \beta(t) \left[ \mathbf{C} \mathbf{I}(t) \circ \frac{\partial \mathbf{S}(t)}{\partial \mathbf{I}(0)} + \mathbf{S}(t) \circ \mathbf{C} \frac{\partial \mathbf{I}(t)}{\partial \mathbf{I}(0)} \right]
\end{align}

where $\circ$ denotes element-wise multiplication with broadcasting: $\mathbf{C} \mathbf{I}(t)$ is broadcast column-wise and $\mathbf{S}(t)$ is broadcast row-wise for element-wise operations with the $2 \times 2$ sensitivity matrices.

\emph{Initial conditions}:
\begin{align}
\frac{\partial \mathbf{S}(0)}{\partial \mathbf{S}(0)} &= \mathbf{I}_2, \quad \frac{\partial \mathbf{I}(0)}{\partial \mathbf{S}(0)} = \mathbf{0} \\
\frac{\partial \mathbf{S}(0)}{\partial \mathbf{I}(0)} &= \mathbf{0}, \quad \frac{\partial \mathbf{I}(0)}{\partial \mathbf{I}(0)} = \mathbf{I}_2
\end{align}
The Fisher information matrix is:
\begin{equation}
\mathbf{J} = \frac{1}{\sigma^2} \mathbf{G}^T \mathbf{G} \in \mathbb{R}^{5 \times 5}
\end{equation}
The diagonal element for connectivity is:
\begin{align}
J_{\theta\theta} &= \frac{1}{\sigma^2} \sum_{t=0}^T \sum_{i=1}^2 \left(\frac{\partial \mu_i(t)}{\partial \theta}\right)^2 \\
&= \frac{1}{\sigma^2} \sum_{t=0}^T \left[ (\beta(t) S_1(t) I_2(t))^2 + (\beta(t) S_2(t) I_1(t))^2 \right] \\
&= \sigma^{-2} \sum_{t=0}^{T-1} \beta(t)^2\left[ S_1^2(t) I_2^2(t) + S_2^2(t) I_1^2(t) \right]
\end{align}

The covariance matrix of any unbiased estimator $\hat{\boldsymbol{\phi}}$ satisfies:
\begin{equation}
\text{Cov}(\hat{\boldsymbol{\phi}}) \succeq \mathbf{J}^{-1}
\end{equation}

To find the exact variance bound for $\theta$ without assuming independence from initial conditions, we partition the Fisher Information Matrix:
\begin{equation}
\mathbf{J} = \begin{bmatrix}
J_{\theta\theta} & \mathbf{j}^T \\
\mathbf{j} & \mathbf{J}_{\text{IC}}
\end{bmatrix}
\end{equation}
where:
\begin{itemize}
\item $J_{\theta\theta} = \frac{1}{\sigma^2} \sum_{t=0}^T \beta(t)^2[ S_1^2(t) I_2^2(t) + S_2^2(t) I_1^2(t) ]$ is the Fisher information for $\theta$
\item $\mathbf{j} = \begin{bmatrix} J_{\theta,S_1(0)} & J_{\theta,I_1(0)} & J_{\theta,S_2(0)} & J_{\theta,I_2(0)} \end{bmatrix}^T$ are cross-correlations
\item $\mathbf{J}_{\text{IC}} \in \mathbb{R}^{4 \times 4}$ is the Fisher information matrix for initial conditions
\end{itemize}

Using the Schur complement, the exact variance bound for $\theta$ is:
\begin{equation}
\text{Var}(\hat{\theta}) \geq [\mathbf{J}^{-1}]_{11} = \frac{1}{J_{\theta\theta} - \mathbf{j}^T \mathbf{J}_{\text{IC}}^{-1} \mathbf{j}}
\end{equation}

The correlation penalty term $\mathbf{j}^T \mathbf{J}_{\text{IC}}^{-1} \mathbf{j} > 0$ reduces the effective Fisher information, making the variance bound larger (worse) than the diagonal approximation $\frac{1}{J_{\theta\theta}}$.

\subsubsection{Cross-Correlation Analysis}

The cross-correlation terms are:
\begin{align}
J_{\theta,S_k(0)} &= \frac{1}{\sigma^2} \sum_{t=0}^{T-1} \sum_{i=1}^2 \frac{\partial \mu_i(t)}{\partial \theta} \frac{\partial \mu_i(t)}{\partial S_k(0)} \\
J_{\theta,I_k(0)} &= \frac{1}{\sigma^2} \sum_{t=0}^{T-1} \sum_{i=1}^2 \frac{\partial \mu_i(t)}{\partial \theta} \frac{\partial \mu_i(t)}{\partial I_k(0)}
\end{align}

\subsubsection{Fisher Information Matrix for Initial Conditions}

The Fisher information matrix for initial conditions $\mathbf{J}_{\text{IC}} \in \mathbb{R}^{4 \times 4}$ has elements:
\begin{equation}
[\mathbf{J}_{\text{IC}}]_{kl} = \frac{1}{\sigma^2} \sum_{t=0}^T \sum_{i=1}^2 \frac{\partial \mu_i(t)}{\partial \phi_k} \frac{\partial \mu_i(t)}{\partial \phi_l}
\end{equation}
where $\phi_k, \phi_l \in \{S_1(0), I_1(0), S_2(0), I_2(0)\}$.

Explicitly, the $4 \times 4$ matrix is:
\begin{equation}
\mathbf{J}_{\text{IC}} = \frac{1}{\sigma^2} \begin{bmatrix}
\sum_{t,i} \left(\frac{\partial \mu_i(t)}{\partial S_1(0)}\right)^2 & \sum_{t,i} \frac{\partial \mu_i(t)}{\partial S_1(0)} \frac{\partial \mu_i(t)}{\partial I_1(0)} & \sum_{t,i} \frac{\partial \mu_i(t)}{\partial S_1(0)} \frac{\partial \mu_i(t)}{\partial S_2(0)} & \sum_{t,i} \frac{\partial \mu_i(t)}{\partial S_1(0)} \frac{\partial \mu_i(t)}{\partial I_2(0)} \\
\sum_{t,i} \frac{\partial \mu_i(t)}{\partial I_1(0)} \frac{\partial \mu_i(t)}{\partial S_1(0)} & \sum_{t,i} \left(\frac{\partial \mu_i(t)}{\partial I_1(0)}\right)^2 & \sum_{t,i} \frac{\partial \mu_i(t)}{\partial I_1(0)} \frac{\partial \mu_i(t)}{\partial S_2(0)} & \sum_{t,i} \frac{\partial \mu_i(t)}{\partial I_1(0)} \frac{\partial \mu_i(t)}{\partial I_2(0)} \\
\sum_{t,i} \frac{\partial \mu_i(t)}{\partial S_2(0)} \frac{\partial \mu_i(t)}{\partial S_1(0)} & \sum_{t,i} \frac{\partial \mu_i(t)}{\partial S_2(0)} \frac{\partial \mu_i(t)}{\partial I_1(0)} & \sum_{t,i} \left(\frac{\partial \mu_i(t)}{\partial S_2(0)}\right)^2 & \sum_{t,i} \frac{\partial \mu_i(t)}{\partial S_2(0)} \frac{\partial \mu_i(t)}{\partial I_2(0)} \\
\sum_{t,i} \frac{\partial \mu_i(t)}{\partial I_2(0)} \frac{\partial \mu_i(t)}{\partial S_1(0)} & \sum_{t,i} \frac{\partial \mu_i(t)}{\partial I_2(0)} \frac{\partial \mu_i(t)}{\partial I_1(0)} & \sum_{t,i} \frac{\partial \mu_i(t)}{\partial I_2(0)} \frac{\partial \mu_i(t)}{\partial S_2(0)} & \sum_{t,i} \left(\frac{\partial \mu_i(t)}{\partial I_2(0)}\right)^2
\end{bmatrix}
\end{equation}

where the summations are over $t \in \{0,1,\ldots,T-1\}$ and $i \in \{1,2\}$.

Each element involves the sensitivity derivatives $\frac{\partial \mu_i(t)}{\partial \phi_k}$ which depend on:
\begin{align}
\frac{\partial \mu_i(t)}{\partial \phi_k} &= \beta(t) \left[ \frac{\partial S_i(t)}{\partial \phi_k} \sum_j C_{ij} I_j(t) + S_i(t) \sum_j C_{ij} \frac{\partial I_j(t)}{\partial \phi_k} \right]
\end{align}

These partial derivatives $\frac{\partial S_i(t)}{\partial \phi_k}$ and $\frac{\partial I_j(t)}{\partial \phi_k}$ evolve according to the sensitivity equations derived earlier.

The matrix $\mathbf{J}_{\text{IC}}$ captures how informative the epidemic data is about each initial condition and their correlations. Regions with stronger epidemic coupling (higher $\theta$) will have more correlated initial condition information.

These correlations arise because:
1. \emph{Initial conditions affect epidemic trajectories}: Changes in $S_i(0), I_i(0)$ propagate through the nonlinear dynamics
2. \emph{Connectivity affects same trajectories}: The parameter $\theta$ also influences how regions interact over time
3. \emph{Observation overlap}: The same infected populations appear in both connectivity and initial condition signals

